% !TEX root = ../notes.tex

\documentclass[../notes.tex]{subfiles}

\begin{document}

\section{October 8}
Today, we continue talking about quasicoherent sheaves.

\subsection{Six Functors for Stacks}
We now pass to stacks.
\begin{definition}[stack]
	A \textit{stack} is an fpqc sheaf from $\mathbb E_\infty$-rings to spaces.
\end{definition}
\begin{definition}
	Fix a stack $X$, and let $\mc U\to X$ be an affine fpqc cover. Then we define $n\mathrm{Pr}_X$ to be the limit of the $n\mathrm{Pr}_U$ as $U$ varies over $U\in\mc U$.
\end{definition}
\begin{remark} \label{rem:npr-for-stack}
	This is distinct from trying to define $n\mathrm{Pr}_X$ inductively from $D(X)$. Intuitively, the reason is that $D(X)$ is probably unable to remember the full content of the $\infty$-stack $X$, but the Stefanich ring should.
\end{remark}
\begin{remark}
	Here is another way to state this definition: the functor $n\mathrm{Pr}_\bullet$ on stacks is defined as the Kan extension of the functor $n\mathrm{Pr}_\bullet\colon\mathbb E_\infty\mathrm{Ring}\opp\to(n+1)\mathrm{Pr}$. This Kan extension takes limits with the underlying $(\infty,1)$-category structure, even though $(n+1)\mathrm{Pr}$ frequently admits module structure allowing it to pretend to be some weirder $(\infty,?)$-category.
\end{remark}
To do descent, we need some countability hypothesis.
\begin{definition}[bicountably presented]
	A morphism $f\colon Y\to X$ of stacks is \textit{bicountably presented} if and only if the following holds: for any pullback square
	% https://q.uiver.app/#q=WzAsNCxbMCwxLCJcXFNwZWMgQSJdLFsxLDEsIlgiXSxbMSwwLCJZIl0sWzAsMCwiXFxTcGVjIEFcXHRpbWVzX1hZIl0sWzIsMSwiZiJdLFszLDAsImYnIiwyXSxbMCwxXSxbMywyXV0=&macro_url=https%3A%2F%2Fraw.githubusercontent.com%2FdFoiler%2Fnotes%2Fmaster%2Fnir.tex
	\[\begin{tikzcd}[cramped]
		{\Spec A\times_XY} & Y \\
		{\Spec A} & X
		\arrow[from=1-1, to=1-2]
		\arrow["{f'}"', from=1-1, to=2-1]
		\arrow["f", from=1-2, to=2-2]
		\arrow[from=2-1, to=2-2]
	\end{tikzcd}\]
	the map $f'$ is a countable colimit of countably presented $A$-algebras. We let $E$ be the subcategory of bicountably presented morphisms.
\end{definition}
\begin{remark}
	Tomer suspects that being countably presented can be detected on $\pi_*$, so these notions admit classical (graded) counterparts.
\end{remark}
The goal of this lecture is to prove the following theorem.
\begin{theorem} \label{thm:6ff-stacks}
	Let $f\colon Y\to X$ be a bicountably presented morphism of stacks. Choose $n\ge1$.
	\begin{listalph}
		\item The functor $f_n^*\colon n\mathrm{Pr}_X\to n\mathrm{Pr}_Y$ admits a left adjoint $f_{n\sharp}$. This map lives in $(n+1)\mathrm{Pr}_X$ and commutes with base-change.
		\item Dualizability: if $n\ge2$, then $f_{n\sharp}(1)$ is dualizable in $n\mathrm{Pr}_X$.
		\item Ambidexterity: if $n\ge2$, then $f_{n\sharp}=f_{n*}$.
	\end{listalph}
\end{theorem}
Intuitively, we are saying that (almost) all maps are \'etale for $n\ge1$, and they are proper for $n\ge2$. For example, for finite \'etale maps of schemes then $f_*$ is a left adjoint of $f^*$. In general, an \'etale map should have $f^*=f^!$, which then has left adjoint $f_!$; there is a similar discussion for proper maps because $f_!=f_*$ for proper maps.
\begin{remark}
	The affine case does not require $n\ge2$, which is perhaps remarkable.
\end{remark}
\begin{proof}[Proof of \Cref{thm:6ff-stacks}]
	Let's start with the case where $X=\Spec A$. We show (a) and (b).
	\begin{listalph}
		\item By hypothesis on $f$, we may write $Y=\colim_I\Spec B_i$ where $I$ is countable, and each map $A\to B_i$ is countably generated. Then
		\[n\mathrm{Pr}_Y=\lim_{I^*}n\mathrm{Pr}_{B_i},\]
		where $I*$ means that we are taking the limit over the pullback morphisms. We already have a left adjoint for morphisms of the affine schemes, so
		\[n\mathrm{Pr}_Y=\colim_{I_\sharp}n\mathrm{Pr}_{B_i}.\]
		There are now natural maps $f_{i,n\sharp}\colon n\mathrm{Pr}_{B_i}\to n\mathrm{Pr}_A$, and they assemble into the desired left adjoint $f_{n\sharp}\colon n\mathrm{Pr}_Y\to n\mathrm{Pr}_A$.
		\item Recall that $I_*=I_\sharp$. Thus, $f_{n\sharp}(1)$ is isomorphic to $\colim_{I_\sharp}f_{n\sharp}(1)$. We now conclude because a countable colimit of dualizable objects in $2\mathrm{Pr}$ along left adjoints is again dualizable. (This argument does not need $n\ge2$.) Observe that there is something tricky happening here: 
	\end{listalph}
	Let's now work globally with general $X$. As before, we can write
	\[(n+1)\mathrm{Pr}_X=\lim_{\Spec A\to X}(n+1)\mathrm{Pr}_A.\]
	Here, the limit may as well be taken in $1\mathrm{Pr}$. For (a), note that each map $n\mathrm{Pr}_Y\to n\mathrm{Pr}_A$ already has the desired left adjoint, and it turns out that compatibility with base-change allows these maps to glue to a map $n\mathrm{Pr}_Y\to n\mathrm{Pr}_X$. For (b), we don't say anything other than asserting that dualizability can be checked locally.\footnote{I am guessing that one can glue together a dual object.}

	It remains to handle ambidexterity. We will use \Cref{prop:get-ambi}. Thus, we want to show that the natural maps ${\id}\to f_n^*f_{n\sharp}$ and $f_{n\sharp}f_n^*\to\id$ have right adjoints.
	\begin{itemize}
		\item For the counit $f_{n\sharp}f_n^*\to\id$, by base-change, it is enough to just exhibit a right adjoint for $f_{n\sharp}(1)\to 1$. By some formula for dualizability, it is equivalent to ask for $1\to f_{n*}1$ to have a left adjoint. This latter map is simply $f_{n-1}^*\colon(n-1)\mathrm{Pr}_X\to(n-1)\mathrm{Pr}_Y$, so we get the result by the existence of $f_{n-1,\sharp}$. Note that we have used $n\ge2$ so that $n-1\ge1$.
		\item For the unit ${\id}\to f_n^*f_{n\sharp}$, we build the following diagonal diagram.
		% https://q.uiver.app/#q=WzAsNSxbMCwwLCJZIl0sWzEsMSwiWVxcdGltZXNfWFkiXSxbMiwxLCJZIl0sWzEsMiwiWSJdLFsyLDIsIlgiXSxbMyw0LCJmIiwyXSxbMiw0LCJmIl0sWzEsMiwicCIsMl0sWzEsMywicSJdLFswLDMsIiIsMSx7ImN1cnZlIjoyLCJsZXZlbCI6Miwic3R5bGUiOnsiaGVhZCI6eyJuYW1lIjoibm9uZSJ9fX1dLFswLDIsIiIsMSx7ImN1cnZlIjotMiwibGV2ZWwiOjIsInN0eWxlIjp7ImhlYWQiOnsibmFtZSI6Im5vbmUifX19XSxbMCwxLCJcXERlbHRhIiwxXV0=&macro_url=https%3A%2F%2Fraw.githubusercontent.com%2FdFoiler%2Fnotes%2Fmaster%2Fnir.tex
		\[\begin{tikzcd}[cramped]
			Y && \\
			& {Y\times_XY} & Y \\
			& Y & X
			\arrow["\Delta"{description}, from=1-1, to=2-2]
			\arrow[curve={height=-12pt}, equals, nfold, from=1-1, to=2-3]
			\arrow[curve={height=12pt}, equals, nfold, from=1-1, to=3-2]
			\arrow["p"', from=2-2, to=2-3]
			\arrow["q", from=2-2, to=3-2]
			\arrow["f", from=2-3, to=3-3]
			\arrow["f"', from=3-2, to=3-3]
		\end{tikzcd}\]
		Let the left downward projection be $p$. Now, ${\id}=p_{n\sharp}\Delta_{n\sharp}\Delta_n^*q_n^*$, and we want a map to $p_{n\sharp}q_n^*=f_n^*f_{n\sharp}$. (This second equality holds by base-change.) The desired map turns to be the map we constructed in the counit point for the morphism $\Delta$!
		\qedhere
	\end{itemize}
\end{proof}
\begin{remark}
	The proof of (c) not only shows that $f_{n\sharp}$ and $f_{n*}$ are the same, but also the identification is compatible with base-change. One could attempt to give an identification by building up from affines, but then it is not obvious that this construction is compatible with base-change.
\end{remark}
\begin{example}
	Another way to state (b) is that $(n-1)\mathrm{Pr}_Y$ is dualizable in $n\mathrm{Pr}_X$.
\end{example}

\subsection{Applications}
Let's apply \Cref{thm:6ff-stacks} to show some things.
\begin{corollary}[K\"unneth]
	Fix bicountably presented morphisms $f\colon Y\to X$ and $f'\colon Y'\to X$. For any $n\ge2$,
	\[(n-1)\mathrm{Pr}_{Y\times_XY'}\cong(n-1)\mathrm{Pr}_Y\otimes(n-1)\mathrm{Pr}_{Y'},\]
	where the isomorphism (and tensor product) takes place in $n\mathrm{Pr}_X$.
\end{corollary}
\begin{proof}
	By \Cref{thm:6ff-stacks}, we have (proper) base-change for these morphisms. Let's see this. We consider the following diagram.
	% https://q.uiver.app/#q=WzAsNCxbMCwwLCJZXFx0aW1lc19YWSciXSxbMSwwLCJZJyJdLFswLDEsIlkiXSxbMSwxLCJYIl0sWzIsMywicCIsMl0sWzEsMywicCciXSxbMCwxLCJmIl0sWzAsMiwiZiciLDJdLFswLDMsImQiLDFdXQ==&macro_url=https%3A%2F%2Fraw.githubusercontent.com%2FdFoiler%2Fnotes%2Fmaster%2Fnir.tex
	\[\begin{tikzcd}[cramped]
		{Y\times_XY'} & {Y'} \\
		Y & X
		\arrow["f", from=1-1, to=1-2]
		\arrow["{f'}"', from=1-1, to=2-1]
		\arrow["d"{description}, from=1-1, to=2-2]
		\arrow["{p'}", from=1-2, to=2-2]
		\arrow["p"', from=2-1, to=2-2]
	\end{tikzcd}\]
	Now, using the projection formula and base-change,
	\begin{align*}
		(n-1)\mathrm{Pr}_{Y\times_XY'} &= d_*1_{Y\times_XY'} \\
		&\cong p_*'f_*(f^*1_{Y'}\otimes {f'}^*1_Y) \\
		&\stackrel*\cong p_*'(1_{Y'}\otimes f_*{f'}^*1_Y) \\
		&\cong p_*'(1_{Y'}\otimes{p'}^*p_*1_Y) \\
		&\cong p_*'(1_{Y'}\otimes p_*1_Y \\
		&= (n-1)\mathrm{Pr}_Y\otimes(n-1)\mathrm{Pr}_{Y'}.
	\end{align*}
	Here, $\stackrel*\cong$ is an application of base-change, which needs the full force of \Cref{thm:6ff-stacks}. Roughly speaking, what is going on is that formally base-change holds for $(-)_\sharp$ via some checks doable with (a), and base-change for our morphisms holds for $(-)_*$ by using \Cref{thm:6ff-stacks}(c).
\end{proof}
\begin{corollary}
	Fix maps $f\colon Y\to X$ and $f'\colon Y'\to X$ of $\aleph_1$-compact stacks. For $n\ge2$, the right adjoints $f_{n*}$ of $f_n^*$ commute with small colimits, tensor products, and pullbacks, and there is a K\"unneth formula.
\end{corollary}
\begin{proof}
	Omitted.
\end{proof}
\begin{remark}
	Without the countability hypothesis, it may happen that $f_{n*}$ is not in $(n+1)\mathrm{Pr}_X$, but it will live in $\op{Mod}_{n\mathrm{Pr}_X}\mathrm{Pr}^L$. Recall that these are different by \Cref{rem:npr-for-stack}.
\end{remark}
\begin{corollary}
	There is a functor from $\aleph_1$-compact stacks to the opposite category of Stefanich rings given by sending $X$ to
	\[(\OO(X),D(X),1\mathrm{Pr}_X,2\mathrm{Pr}_X,\ldots),\]
	and it commutes with countable colimits and finite limits.
\end{corollary}
\begin{proof}
	The construction allows us to see that we have built a Stefanich ring, and it commutes with colimits by construction of the $n\mathrm{Pr}_X$s. Thus, we will focus our energies on showing that it commutes with finite limits. It does send the final object $\Spec\mathbb S$ (which is the stack associated to $\mathbb S$) to the final opposite Stefanich ring $\mathbb S$.

	It remains to show that our functor preserves fiber products. Roughly speaking, one expects the fiber product of $Y\to X$ and $Y'\to X$ to be a Stefanich ring which looks like
	\[(\OO(Y)\otimes_{\OO(X)}\OO(Y'),D(Y)\otimes_{D(X)}D(Y'),1\mathrm{Pr}_Y\otimes_{1\mathrm{Pr}_X}1\mathrm{Pr}_{Y'},\ldots),\]
	but this is not actually a Stefanich ring. However, there is a process by which one can force this to be a Stefanich ring, and it turns out to compare this to
	\[\left(\OO(Y)\otimes^{D(X)}\OO(Y'),D(Y)\otimes^{1\mathrm{Pr}_X}D(Y'),1\mathrm{Pr}_Y\otimes^{2\mathrm{Pr}_X}1\mathrm{Pr}_{Y'},\ldots\right),\]
	which is also not literally a Stefanich ring (here the superscripts mean that the tensor products take place in that category), but it agrees with the correct fiber product Stefanich ring.
\end{proof}

\end{document}